% =====================================================================
% COSC 3337 — Homework 4 — Storage and Normalization
% =====================================================================
\documentclass[11pt]{article}

\usepackage[margin=1in]{geometry}
\usepackage{amsmath,amssymb}
\usepackage{booktabs}
\usepackage{enumitem}
\usepackage{fancyhdr}
\usepackage{xcolor}
\usepackage{textcomp}
\usepackage{listings}
\usepackage[most]{tcolorbox}
\usepackage{array}
\usepackage{tabularx}
\usepackage{tikz}
\usepackage{needspace}
\usetikzlibrary{positioning,arrows.meta}
\usepackage[hidelinks]{hyperref}

% ---- edit these ------------------------------------------------------
\newcommand{\hwnum}{4}
\newcommand{\released}{TBD}
\newcommand{\due}{Sunday, November 22, 2026, 11:59 pm}
% ---------------------------------------------------------------------

\definecolor{boxred}{HTML}{C9605A}
\definecolor{boxredbg}{HTML}{FBF3F2}
\definecolor{boxblue}{HTML}{4A6F96}
\definecolor{boxbluebg}{HTML}{F2F6FA}
\definecolor{codegray}{HTML}{6B6B6B}

\newtcolorbox{warnbox}[1]{enhanced, breakable, colback=boxredbg, colframe=boxred,
  fonttitle=\bfseries\small, title={#1}, boxrule=0.6pt, arc=1pt,
  left=6pt, right=6pt, top=4pt, bottom=4pt}
\newtcolorbox{notebox}[1]{enhanced, breakable, colback=boxbluebg, colframe=boxblue,
  fonttitle=\bfseries\small, title={#1}, boxrule=0.6pt, arc=1pt,
  left=6pt, right=6pt, top=4pt, bottom=4pt}

\lstset{
  basicstyle=\ttfamily\small,
  columns=fullflexible,
  keepspaces=true,
  upquote=true,
  showstringspaces=false,
  frame=leftline,
  framerule=0.8pt,
  rulecolor=\color{codegray},
  xleftmargin=1.2em,
  framexleftmargin=0.6em,
  aboveskip=6pt,
  belowskip=6pt,
  commentstyle=\color{codegray},
  breaklines=true,
}
\lstdefinestyle{py}{language=Python}
\lstdefinestyle{sql}{language=SQL, morekeywords={LIMIT,COALESCE,EXCEPT,WITH,OVER,PARTITION}}
\newcommand{\code}[1]{\texttt{#1}}
\newcommand{\pts}[1]{\hfill\textit{(#1 points)}}
\newcommand{\pt}[1]{\textit{(#1~pt\ifnum#1>1 s\fi)}}

\pagestyle{fancy}
\fancyhf{}
\lhead{\small COSC 3337 --- Homework \hwnum}
\rhead{\small Page \thepage}
\renewcommand{\headrulewidth}{0.4pt}

\setlength{\parindent}{0pt}
\setlength{\parskip}{6pt}
\setlength{\emergencystretch}{3em}
\setlist{itemsep=2pt, topsep=4pt}

\begin{document}

% =====================================================================
\begin{center}
  {\LARGE\bfseries Homework \hwnum: Storage and Normalization}\\[4pt]
  {\large COSC 3337 --- Databases --- Fall 2026}\\[6pt]
  \begin{tabular}{@{}ll@{}}
    \textbf{Released} & \released \\
    \textbf{Due}      & \due \\
    \textbf{Total}    & 100 points, plus 5 bonus
  \end{tabular}
\end{center}
\vspace{2pt}\hrule\vspace{10pt}

\section*{Before you start}

\textbf{The scenario.} Soundwave has bought \emph{Soundwave Live}, a small
concert-ticketing startup. Soundwave Live ran its whole business out of one
spreadsheet, and you have been handed its export: one MySQL table,
\code{ticket\_sale\_wide}, in a new database, \code{soundwave\_live}. Each row
is one \emph{line} of a ticket order: an order is for one concert, and can buy
tickets in several seat sections.

\begin{center}
\small
\begin{tabularx}{\linewidth}{@{}l X@{}}
\toprule
\textbf{Columns of \code{ticket\_sale\_wide}} & \textbf{Meaning}\\
\midrule
\code{order\_id}, \code{order\_date} & the order, and when it was placed. Order ids increase over time.\\
\code{listener\_id}, \code{username}, \code{email} & the Soundwave listener who placed it\\
\code{concert\_id}, \code{concert\_date} & the concert the order is for\\
\code{artist\_id}, \code{artist\_name} & who is playing\\
\code{venue\_id}, \code{venue\_name}, \code{city}, \code{capacity} & where\\
\code{seat\_section}, \code{price} & a seat section (\code{GA}, \code{FLOOR}, \code{VIP}, \ldots) and its price per ticket at this concert\\
\code{quantity} & how many tickets in that section this order bought\\
\bottomrule
\end{tabularx}
\end{center}

Soundwave Live's staff confirmed these \textbf{business rules}. They are the
set of functional dependencies $F$ for the whole homework:

\begin{center}
\begin{tabular}{@{}ll@{}}
\toprule
\textbf{R1} & \code{listener\_id} $\rightarrow$ \code{username}, \code{email}\\
\textbf{R2} & \code{artist\_id} $\rightarrow$ \code{artist\_name}\\
\textbf{R3} & \code{venue\_id} $\rightarrow$ \code{venue\_name}, \code{city}, \code{capacity}\\
\textbf{R4} & \code{concert\_id} $\rightarrow$ \code{concert\_date}, \code{artist\_id}, \code{venue\_id}\\
\textbf{R5} & \code{concert\_id}, \code{seat\_section} $\rightarrow$ \code{price}\\
\textbf{R6} & \code{order\_id} $\rightarrow$ \code{order\_date}, \code{listener\_id}, \code{concert\_id}\\
\textbf{R7} & \code{order\_id}, \code{seat\_section} $\rightarrow$ \code{quantity}\\
\bottomrule
\end{tabular}
\end{center}

\textbf{Part 1} is written (50 points). \textbf{Part 2} is SQL (50 points).
Questions 1.4 and 1.5 analyze \code{ticket\_sale\_wide} on paper, and Part~2 then
does the same work in MySQL, so do 1.5 before 2.3.
\\


\textbf{What to submit:} one zip file containing exactly:
\begin{enumerate}[nosep]
  \item \code{hw4.pdf} --- your answers to Part 1.
  \item \code{fd\_checks.sql} --- your answer to Question 2.1.
  \item \code{fix.sql} --- your answer to Question 2.2.
  \item \code{normalize.sql} --- your answer to Question 2.3.
  \item \code{lossless\_check.sql} --- your answer to Question 2.4.
\end{enumerate}
Do not rename the four \code{.sql} files. The autograder looks for them by name.

\textbf{Typesetting bonus (5 points).}
\begin{itemize}
  \item \textbf{By hand.} You may hand-write Part 1. Submit a PDF made with a
  scanning app, not a photograph, with every question clearly labeled.
  \item \textbf{With \LaTeX\ (5 bonus points).} If you typeset Part 1 in \LaTeX,
  you earn 5 bonus points. A template, \code{hw4\_template.tex}, is in the
  \code{HW4} folder. Submit the compiled PDF, not the \code{.tex} file.
\end{itemize}

\textbf{Do Part 0 first, and early,} so we have time to fix any setup problems.

% =====================================================================
\section*{Part 0 --- Setup \hfill \textit{(0 points, required)}}

\subsection*{0.1 Get the starter files}

On your VM, go to your copy of the course repository, pull the new homework
from \code{upstream}, and load the new database:
\begin{lstlisting}
cd COSC3337-FA26-homeworks
git pull upstream main
cd HW4
bash setup_db.sh
\end{lstlisting}

\code{setup\_db.sh} loads \code{soundwave\_live}, installs the Python packages
the tests need, and checks that everything connects. It is safe to run again at
any time and restores \code{ticket\_sale\_wide} to the original export, so it
is your repair command. It also \emph{deletes} any tables you created in
\code{soundwave\_live}, but never your \code{.sql} files.

\subsection*{0.2 Run the tests}

\begin{lstlisting}
python3 -m pytest tests/ -v
\end{lstlisting}

Every test should run and \emph{fail} with a message about your answer (the
starter files contain placeholders). That means your setup works. Any other
error (cannot connect, access denied, module not found) is a setup problem.
Re-run \code{setup\_db.sh}, and if that does not fix it, email me or come to
office hours.

\subsection*{0.3 Back up your work}

Commit and push to your private \code{cosc3337-work} repository every time you
stop working. From inside \code{HW4}:
\begin{lstlisting}
git add *.sql
git commit -m "<what you did this session>"
git push origin main
\end{lstlisting}
Your VM has no backup.

% =====================================================================
\newpage
\section*{Part 1 --- Written \hfill \textit{(50 points)}}

Show your work for every calculation.

% ---------------------------------------------------------------------
\subsection*{1.1 RAID \pts{11}}

Soundwave Live is buying a server with four 4~TB disks.

\begin{enumerate}[label=(\alph*)]
  \item \pt{3} For RAID 0 (block striping across all four disks), RAID 1 (two
  mirrored pairs, with blocks striped across the pairs) and RAID 5
  (block-interleaved distributed parity), give the usable capacity and the
  number of disk failures the layout survives \emph{no matter which} disks fail.

  \item \pt{4} Under RAID 5, one stripe holds three data blocks on disks 0, 1
  and 2 and their parity block on disk 3. With blocks shortened to 8 bits:
  \[
    D_0 = \texttt{10110010} \qquad D_1 = \texttt{01101100} \qquad D_2 = \texttt{11100001}
  \]
  \begin{enumerate}[label=\arabic*., nosep]
    \item Compute the parity block $P = D_0 \oplus D_1 \oplus D_2$, showing each step.
    \item Disk 1 fails. Show how the system rebuilds $D_1$ from the blocks that
    survive, and check that you got $D_1$ back.
  \end{enumerate}

  \item \pt{4} Soundwave Live will store two things on this server: the
  \emph{play log}, which receives a small write every time anyone plays a song,
  and the \emph{ticket archive}, which is written once a night and read by
  reports all day. Choose RAID 1 or RAID 5 for each, and justify both choices.
\end{enumerate}

% ---------------------------------------------------------------------
\subsection*{1.2 Slotted pages \pts{7}}

\code{ticket\_sale\_wide} has variable-length records (\code{email} and
\code{venue\_name} vary in length), so a DBMS would store them in slotted
pages. A small 240-byte block holds four of its records in a slotted page. The
header takes 4 bytes plus 4 bytes per slot, and records are packed at the end
of the block. Right now the end of free space is 130, leaving 110 free bytes.

\begin{center}
\small
\begin{tabular}{@{}lccc@{}}
\toprule
\textbf{Slot} & \textbf{(offset, size) now} & \textbf{After (a)} & \textbf{After (b)}\\
\midrule
1 & (215, 25) & & \\
2 & (175, 40) & & \\
3 & (160, 15) & & \\
4 & (130, 30) & & \\
5 & --- & & \\
\midrule
End of free space & 130 & & \\
Free bytes & 110 & & \\
\bottomrule
\end{tabular}
\end{center}

\begin{enumerate}[label=(\alph*)]
  \item \pt{4} Delete the record in slot 1 and compact the page. Fill in the
  header.

  \item \pt{3} Insert a 50-byte record into a new slot 5. Fill in the header.
\end{enumerate}

% ---------------------------------------------------------------------
\subsection*{1.3 File organization \pts{7}}

For each workload, choose a \emph{heap} file, a \emph{sequential} file or
\emph{columnar} storage, and justify your choice.

\begin{enumerate}[label=(\alph*)]
  \item \pt{2} Soundwave's play log: millions of rows appended every day, read
  only by a nightly job that scans all of it.

  \item \pt{3} Ticket orders, stored by \code{order\_id}. Invoicing reads
  them in ranges (``orders 101000 to 101999''), and new orders almost always
  arrive in increasing \code{order\_id}. Also say what happens when a late
  order with a \emph{smaller} \code{order\_id} arrives and the block it
  belongs in is full.

  \item \pt{2} An analytics query computes the average \code{price} for each
  \code{seat\_section} over 500 million rows of a 16-column table like
  \code{ticket\_sale\_wide}. If all columns are the same width, roughly what
  fraction of a row layout's bytes does your choice read?
\end{enumerate}

% ---------------------------------------------------------------------
\subsection*{1.4 Functional dependencies in the export \pts{9}}

Here are four rows of \code{ticket\_sale\_wide} (some columns omitted):

\begin{center}
\footnotesize\setlength{\tabcolsep}{4pt}
\begin{tabular}{@{}ccllcrc@{}}
\toprule
\code{order\_id} & \code{listener\_id} & \code{email} & \code{concert\_id} & \code{seat\_section} & \code{price} & \code{quantity}\\
\midrule
100856 & L0071 & \code{silvermaple72\_9@example.com} & C002 & FLOOR & 89.50 & 1\\
100856 & L0071 & \code{silvermaple72\_9@example.com} & C002 & UPPER & 63.00 & 1\\
101276 & L0071 & \code{silvermaple72@example.com}    & C021 & UPPER & 46.50 & 2\\
101386 & L0071 & \code{silvermaple72@example.com}    & C046 & LOWER & 60.50 & 4\\
\bottomrule
\end{tabular}
\end{center}

\begin{enumerate}[label=(\alph*)]
  \item \pt{2} These rows already show the damage done by one kind of
  anomaly. Which kind, and how did it most likely happen? 

  \item \pt{4} Using $F$, compute $\{\code{order\_id}\}^+$ and
  $\{\code{order\_id}, \code{seat\_section}\}^+$ with the attribute closure
  algorithm. Show which rule you apply at each step. What do the results tell
  you about each attribute set?

  \item \pt{3} Show that $\{\code{order\_id}, \code{seat\_section}\}$ is the
  \emph{only} candidate key of \code{ticket\_sale\_wide}.
\end{enumerate}

% ---------------------------------------------------------------------
\subsection*{1.5 BCNF decomposition \pts{16}}

\begin{enumerate}[label=(\alph*)]
  \item \pt{3} Which of R1--R7 violate BCNF on \code{ticket\_sale\_wide}?
  Justify with your answers to~1.4.

  \item \pt{10} Decompose \code{ticket\_sale\_wide} into BCNF with the BCNF
  decomposition algorithm. For each step, give the relation you split, the
  violating dependency $\alpha \rightarrow \beta$, and the two relations it
  produces, $(\alpha \cup \beta)$ and $(R - (\beta - \alpha))$. Finish with
  the list of final relations, with each one's key underlined.

  At each step, split on a rule from R1--R7 that lies entirely inside the
  relation and violates BCNF there, choosing one whose right side contains no
  attribute from the left side of another rule still inside the relation. A
  relation is finished when none of R1--R7 violates BCNF in it.

  \item \pt{3} Choose one step of your decomposition and show, with the test
  from lecture, that it is lossless.
\end{enumerate}

% =====================================================================
\newpage
\section*{Part 2 --- SQL \hfill \textit{(50 points)}}

\subsection*{The files}

\begin{center}
\begin{tabularx}{\linewidth}{@{}l X l@{}}
\toprule
\textbf{File} & \textbf{What it is} & \textbf{You\ldots}\\
\midrule
\code{fd\_checks.sql}      & 2.1: five queries                         & \textbf{edit and submit}\\
\code{fix.sql}             & 2.2: statements that repair the export    & \textbf{edit and submit}\\
\code{normalize.sql}       & 2.3: statements that build your BCNF tables & \textbf{edit and submit}\\
\code{lossless\_check.sql} & 2.4: one query                            & \textbf{edit and submit}\\
\code{soundwave\_live.sql} & the export, loaded by \code{setup\_db.sh}  & don't edit\\
\code{tests/}              & the tests                                 & run\\
\bottomrule
\end{tabularx}
\end{center}

Run a script against the database with, for example,
\begin{lstlisting}
mysql soundwave_live < fix.sql
\end{lstlisting}
\textbf{Do not write \code{USE}} in any file, and do not create other
databases: the autograder runs your files against its own copy of
\code{soundwave\_live}.

\begin{notebox}{How grading works}
The autograder runs the same tests as \code{tests/}, on a \textbf{different
copy of the data}: different listeners, concerts, orders, and different broken
rows. Answers copied from your copy of the data will fail. Each line of the
tables below is graded on its own, except where a line says otherwise. Each of
2.1--2.3 starts from freshly loaded data, so a mistake in one of them does not
cost you points in another. 2.4 uses the tables your \code{normalize.sql}
builds.

The tests reload \code{soundwave\_live} as they run. When they finish, it is
restored to the original export, so any tables you created are gone. Re-run
your scripts if you want them back.
\end{notebox}

% ---------------------------------------------------------------------
\subsection*{2.1 Find the rule violations \pts{15}}

A spreadsheet cannot enforce a functional dependency, so before trusting the
export, you check it. In \code{fd\_checks.sql}, write one query under each
header. Each query finds the \emph{violations} of one rule: the left-hand-side
values that appear in \code{ticket\_sale\_wide} with two or more different
combinations of right-hand-side values.

\begin{center}
\small
\begin{tabular}{@{}lll@{}}
\toprule
\textbf{Query} & \textbf{Checks} & \textbf{Returns exactly these columns, in this order and sorted by them}\\
\midrule
Q1 & R1 & \code{listener\_id}\\
Q2 & R3 & \code{venue\_id}\\
Q3 & R4 & \code{concert\_id}\\
Q4 & R5 & \code{concert\_id}, \code{seat\_section}\\
Q5 & R6 & \code{order\_id}\\
\bottomrule
\end{tabular}
\end{center}

Each violating value appears once. A value that always appears with the same
right-hand side does not appear at all. Keep the
\code{-- Q1} \ldots\ \code{-- Q5} header lines exactly as they are: the
autograder splits the file on them.

The tests run your queries on your copy of the export and on a second copy in
which all five rules are broken.

\needspace{5\baselineskip}
\begin{center}
\small
\begin{tabularx}{\linewidth}{@{}X r@{}}
\toprule
\textbf{What is tested} & \textbf{Points}\\
\midrule
Each of Q1--Q5 returns exactly the right rows, on both copies & $5 \times 3$\\
\bottomrule
\end{tabularx}
\end{center}

% ---------------------------------------------------------------------
\subsection*{2.2 Repair the export \pts{10}}

Write \code{fix.sql}: one or more \code{UPDATE} statements that repair
\code{ticket\_sale\_wide} in place, so that every business rule holds. Use
your queries from 2.1 to find out what needs repairing.

When a left-hand-side value appears with more than one right-hand side, the
correct right-hand side is the newest one: the one on that value's row with the
largest \code{order\_id}. Change nothing else. The autograder's copy is broken
in the same rules as yours, but in different rows, so your statements must work
for \emph{any} value, not just the ones you found.

\begin{notebox}{A MySQL restriction}
MySQL refuses an \code{UPDATE} whose \code{WHERE} or \code{SET} clause has a
subquery that reads the table being updated (error 1093, \emph{``You can't
specify target table \ldots\ for update in FROM clause''}). A subquery
joined to the table in the \code{UPDATE} itself is allowed.
\end{notebox}

Re-run \code{setup\_db.sh} (or \code{mysql < soundwave\_live.sql}) before each
attempt, so you always start from the original export. Afterwards, every query
in \code{fd\_checks.sql} should return no rows.

\begin{center}
\small
\begin{tabularx}{\linewidth}{@{}X r@{}}
\toprule
\textbf{What is tested} & \textbf{Points}\\
\midrule
\code{fix.sql} runs, and afterwards every rule R1--R7 holds & 4\\
Every repaired value is the newest one, and nothing else changed: the table matches the correctly repaired export exactly & 4\\
No rows were added or deleted (\emph{only awarded if every rule holds}) & 2\\
\bottomrule
\end{tabularx}
\end{center}

% ---------------------------------------------------------------------
\subsection*{2.3 Normalize \pts{20}}

Write \code{normalize.sql}, which turns your answer to 1.5 into real tables:
it creates your BCNF relations in \code{soundwave\_live} and fills them from
the repaired \code{ticket\_sale\_wide}. The autograder runs it on a
\emph{correctly} repaired export, so 2.3 does not depend on your answer to 2.2.

\begin{warnbox}{Rules for normalize.sql (checked by the autograder)}
\begin{itemize}[nosep]
  \item Name your tables anything except \code{ticket\_sale\_wide}. Every base
  table in \code{soundwave\_live} other than \code{ticket\_sale\_wide} counts
  as part of your design, so do not leave scratch tables behind.
  \item Use \emph{only} the 16 column names of \code{ticket\_sale\_wide}, with
  the same types: no new columns (no surrogate ids), no renaming. Every one of
  the 16 must appear in at least one table.
  \item Every table has a \code{PRIMARY KEY}.
  \item When a table contains all the columns of another table's primary key,
  and those columns are not exactly its own primary key, declare a
  \code{FOREIGN KEY} on those columns referencing the other table.
  \item Fill every table from \code{ticket\_sale\_wide} with
  \code{INSERT \ldots\ SELECT}. Do not change or drop \code{ticket\_sale\_wide}.
  \item Start with \code{DROP TABLE IF EXISTS} for your tables, so the file can
  run again and again. Drop child tables before the tables they reference.
\end{itemize}
\end{warnbox}

The checks use $F$ (R1--R7). In (b), a table \emph{stores rule R$i$
redundantly} if it contains R$i$'s left side and at least one attribute of its
right side, but R$i$'s left side is not a superkey of the table.

\begin{center}
\small
\begin{tabularx}{\linewidth}{@{}X r@{}}
\toprule
\textbf{What is tested} & \textbf{Points}\\
\midrule
(a) \code{normalize.sql} runs; the tables use exactly the 16 columns; \code{ticket\_sale\_wide} is unchanged & 2\\
(b) For each of R1--R6, no table stores it redundantly & $6 \times 1$\\
\midrule
\multicolumn{2}{@{}l}{\emph{(c)--(f) are only awarded if at least 4 of the 6 checks in (b) pass.}}\\
(c) Every table is in BCNF & 2\\
(d) Joining all your tables gives back exactly the rows of \code{ticket\_sale\_wide} & 6\\
(e) Every table's primary key is a candidate key of that table & 2\\
(f) Every foreign key required by the rules above is declared & 2\\
\bottomrule
\end{tabularx}
\end{center}

% ---------------------------------------------------------------------
\subsection*{2.4 Prove it is lossless \pts{5}}

In 1.5 you argued that one step of your decomposition is lossless. Now check
the whole decomposition against the data. Write \code{lossless\_check.sql}:
one query (it may start with \code{WITH}) that returns
\begin{itemize}[nosep]
  \item every row of \code{ticket\_sale\_wide} that joining your tables does
  not produce, and
  \item every row that joining your tables produces that is not in
  \code{ticket\_sale\_wide}.
\end{itemize}
When your decomposition is lossless, it returns no rows. The columns it
returns are up to you. The MySQL on your VM supports \code{EXCEPT}.

To test it, the autograder builds your tables with \emph{your}
\code{normalize.sql} and then damages them.

\begin{center}
\small
\begin{tabularx}{\linewidth}{@{}X r@{}}
\toprule
\textbf{What is tested} & \textbf{Points}\\
\midrule
Returns at least one row after one row is deleted from your largest table & 2\\
Returns at least one row after one value outside a primary key is changed & 2\\
Returns no rows for your tables as \code{normalize.sql} built them (\emph{only awarded if the first test passes}) & 1\\
\bottomrule
\end{tabularx}
\end{center}

\subsection*{Before you submit}

\begin{enumerate}[nosep]
  \item \code{python3 -m pytest tests/ -v} passes.
  \item Starting from a fresh \code{setup\_db.sh}, your files run without
  errors in this order: \code{fix.sql}, \code{normalize.sql},
  \code{lossless\_check.sql}.
  \item Your zip contains exactly \code{hw4.pdf} and the four \code{.sql}
  files, with their original names.
\end{enumerate}

\end{document}
